\subsection{深度与 Koszul 复形}

设 A 为环，M 为 A-模，\(\mathbf{x} = (x_i)_{i \in I}\) 为 A 中元素的一个族。用 \(u : \mathrm{A}^{(1)} \to \mathrm{A}\) 表示使 \(u(c_i) = x_i\) 对每个 \(i \in I\) 成立的线性形式，并以 \(\mathbf{K}^{\bullet}(\mathbf{x}, \mathrm{M})\) 表示与 u 相伴的复形 \(\mathbf{K}_{\mathrm{A}}^{\bullet}(u, \mathrm{M})\) (A, X, 第 147 页)。\(\mathbf{K}^p(\mathbf{x}, \mathrm{M}) = 0\) 对 p \textless{} 0 成立；对 \(p \geqslant 0\) ，A-模 \(\mathbf{K}^p(\mathbf{x}, \mathrm{M}) = \operatorname{Hom}_{\Lambda} (\boldsymbol{\Lambda}^p(\mathrm{A}^{(1)}), \mathrm{M})\) 典范地等同于 A-模 \(C_I^p(\mathrm{M})\) ，后者由 \(I^p\) 到 M 的交错映射构成 (A, X, 第 153 页)，而微分 \(\partial^p : \mathbf{K}^p(\mathbf{x}, \mathrm{M}) \longrightarrow \mathbf{K}^{p+1}(\mathbf{x}, \mathrm{M})\) 由下述公式给出

\[
\left(\partial^ {p} \boldsymbol {m}\right) \left(\alpha_ {1}, \dots , \alpha_ {p + 1}\right) = \sum_ {j = 1} ^ {p + 1} (- 1) ^ {j + 1} x _ {\alpha_ {j}} \boldsymbol {m} \left(\alpha_ {1}, \dots , \alpha_ {j - 1}, \alpha_ {j + 1}, \dots , \alpha_ {p + 1}\right)
\]

对 \(\boldsymbol{m} \in \mathbf{K}^{p}(\mathbf{x}, \mathrm{M})\) 及 \((\alpha_{1}, \ldots, \alpha_{p+1}) \in \mathrm{I}^{p+1}\) 成立 (A, X, 第 154 页, 公式 (12))。特别地，由此得出复形 \(\mathbf{K}^{\bullet}(\mathbf{x}, \mathrm{M})\) 仅依赖于 M 作为 Z-模的结构以及自同态 \((x_{i})_{\mathrm{M}}\) 。

以 \(\mathrm{H}^{\bullet}(\mathbf{x},\mathrm{M})\) 表示复形 \(\mathbf{K}^{\bullet}(\mathbf{x},\mathrm{M})\) 的同调。A-模 \(\mathrm{H}^{0}(\mathbf{x},\mathrm{M})\) 等同于 \(\mathrm{Hom}_{\Lambda}(\mathrm{A}/\mathrm{J},\mathrm{M})\) ，其中 J 为由 \(x_{i}\) 生成的 A 的理想 (A, X, 第 147 页, 引理 1)。

设 \((\mathrm{M}_{\alpha})_{\alpha\in\mathrm{K}}\) 为一族 A-模，M 为其乘积；复形 \(\mathbf{K}^{\bullet}(\mathbf{x},\mathrm{M})\) 典范同构于诸 \(\mathbf{K}^{\bullet}(\mathbf{x},\mathrm{M}_{\alpha})\) 的乘积复形，因而对每个整数 s，A-模 \(\mathrm{H}^{s}(\mathbf{x},\mathrm{M})\) 等同于诸 \(\mathrm{H}^{s}(\mathbf{x},\mathrm{M}_{\alpha})\) 的乘积 (A, X, 第 28 页, 命题 1)。

定理 1.-- 设 A 为环，J 为 A 的理想，\(\mathbf{x} = (x_i)_{i \in I}\) 为 J 的一个生成族，M 为 A-模。M 关于 J 的深度是 \(\mathbf{N} \cup \{+\infty\}\) 中由满足下述条件的整数 \(n\) 组成的集合的下确界：\(\Pi^n(\mathbf{x}, \mathbf{M}) \neq 0\) 。

令 \(p = \text{prof}_{\Lambda}(\text{J}; \text{M})\) 。考察复形 \(\mathbf{K}^{\bullet}(\mathbf{x}, \text{M})\) 。它的同调被 J 零化 ( \(\Lambda\) , X, 第 148 页, 推论 2)，且诸模 \(\mathbf{K}^{i}(\mathbf{x}, \text{M})\) 中每一个关于 J 的深度等于 \(p\) 或 \(+\infty\) （第 1 节, 注 4）。于是由第 2 节的推论 1 得出：\(\text{H}^{i}(\text{x}, \text{M}) = 0\) 对 \(i < p\) 成立。剩下要证明 \(\text{H}^{p}(\text{x}, \text{M})\) 在 \(p < +\infty\) 时非零。

情形 \(p = 0\) 是显然的，故设 \(0 < p < +\infty\) 且 \(\mathrm{H}^p(\mathbf{x},\mathrm{M}) = 0\) 。设 L 为 A-模 \(\mathrm{A / J}\) 的一个自由分解；以 C 表示复形 \(\mathrm{Homgr}_{\mathrm{A}}(\mathrm{L},\mathrm{M})\) 。\(\Lambda\) -模 \(\mathrm{H}^i(\mathrm{C})\) 同构于 \(\mathrm{Ext}_{\mathrm{A}}^i(\mathrm{A / J},\mathrm{M})\) (A, X, 第 100 页, 定理 1)；因而它当 i \textless{} p 时为零。~于是，对 i \textless{} p 有典范正合列

\[
0 \rightarrow \mathrm{B} ^ {i} (\mathrm{C}) \rightarrow \mathrm{C} ^ {i} \rightarrow \mathrm{B} ^ {i + 1} (\mathrm{C}) \rightarrow 0.
\]

A-模 \(C^{i}\) 是一些同构于 M 的 A-模的乘积；故 \(\mathrm{H}^{s}(\mathbf{x},\mathrm{C}^{i})=0\) 对 \(s\leqslant p\) 成立。由上述正合列及 A, X, 第 150 页 推出：连接同态 \(\partial^{s}:\mathrm{H}^{s}(\mathbf{x},\mathrm{B}^{i+1}(\mathrm{C}))\longrightarrow\mathrm{H}^{s+1}(\mathbf{x},\mathrm{B}^{i}(\mathrm{C}))\) 对 \(s\leqslant p\) 及 i\textless p 是单射；由于 \(\mathrm{B}^{0}(\mathrm{C})=0\) ，由此得出 \(\mathrm{H}^{p-i}(\mathbf{x},\mathrm{B}^{i+1}(\mathrm{C}))\) 当 i\textless p 时为零。~特别地，有 \(\mathrm{II}^{1}(\mathbf{x},\mathrm{B}^{p}(\mathrm{C}))=0\) ，于是正合列

\[
0 \to \mathrm{B} ^ {p} (\mathrm{C}) \to \mathrm{Z} ^ {p} (\mathrm{C}) \to \mathrm{H} ^ {p} (\mathrm{C}) \to 0
\]

给出一个满射 \(\mathrm{H}^{0}(\mathbf{x},\mathrm{Z}^{p}(\mathrm{C}))\longrightarrow\mathrm{H}^{0}(\mathbf{x},\mathrm{H}^{p}(\mathrm{C}))\) 。由于 \(\mathrm{H}^{p}(\mathrm{C})\) 同构于 \(\operatorname{Ext}_{\mathrm{A}}^{p}(\mathrm{A}/\mathrm{J},\mathrm{M})\) ，而后者非零且被 J 零化，故有 \(\mathrm{H}^{0}(\mathbf{x},\mathrm{H}^{p}(\mathrm{C}))\neq0\) ，从而 \(\mathrm{H}^{0}(\mathbf{x},\mathrm{Z}^{p}(\mathrm{C}))\neq0\) ，于是 \(\mathrm{H}^{0}(\mathbf{x},\mathrm{C}^{p})\neq0\) 。但这蕴含 \(\mathrm{H}^{0}(\mathbf{x},\mathrm{M})\neq0\) ，与假设矛盾。因此 \(\mathrm{H}^{p}(\mathbf{x},\mathrm{M})\neq0\) ，证明完毕。

推论 1.--- 设理想 J 有限生成且 JM ≠ M。则 prof \(_{A}\) (J; M) 有限且 ≤ Card(I)；为使它等于 Card(I)，必须且只需族 x 对 M 是完全割的 (A, X, 第 157 页, 定义 2)。

首先设 I 有限，并以 \(r\) 表示其基数。A-模 \(\mathrm{H}^r (\mathbf{x},\mathrm{M})\) 典范地同构于 \(\mathrm{H}_0(\mathbf{x},\mathrm{M})\) ，后者又同构于 M/JM (A, X, 第 155 页)；因而不等式 \(\mathrm{prof}_{\Lambda}(\mathrm{J};\mathrm{M})\leqslant r\) 由定理 1 得出。为使等号成立，必须且只需 A-模 \(\mathrm{H}^i (\mathbf{x},\mathrm{M})\) 对 \(i < r\) 均为零，而这恰好意味着族 \(\mathbf{x}\) 对 M 是完全割的 (A, X, 第 157 页)。

根据前文，\(\mathrm{prof}_{\mathrm{A}}(\mathrm{J};\mathrm{M})\) 有限；剩下要证明：若族 \(\mathbf{x}\) 对 M 完全割，则集 I 有限。但条件 \(\mathrm{H}_{1}(\mathbf{x},\mathrm{M}) = 0\) (A, X, 第 157 页, 定义 2) 蕴含我们有正合列

\[
\boldsymbol {\Lambda} ^ {2} (\mathrm{A} ^ {(\mathrm{I})}) \otimes_ {\mathrm{A}} \mathrm{M} \xrightarrow {\partial_ {2}} \mathrm{M} ^ {(\mathrm{I})} \xrightarrow {\partial_ {1}} \mathrm{JM} \to 0,
\]

其中 \(\partial_{2}\) 的像含于 \(\mathrm{JM}^{(1)}\) 。用 A/J 作张量积，我们导出从 \((\mathrm{M} / \mathrm{JM})^{(1)}\) 到 \(\mathrm{JM} / \mathrm{J}^2\mathrm{M}\) 上的一个 A-线性同构。而后一模是有限型的，因为 J 与 M 都是；由于 M/JM 非零，于是得出集 I 有限。

推论 2. 设 A 为局部环，J 为 A 的异于 A 的有限生成理想，M 为非零的有限生成 A-模。令 \(r = [\mathrm{J} / \mathfrak{m}_{\mathrm{A}}\mathrm{J} : \kappa_{\mathrm{A}}]\) 。有 \(\operatorname{prof}_{\Lambda}(\mathrm{J};\mathrm{M}) \leqslant r\) ；当且仅当 J 由一个对 M 完全割的族生成时等号成立。这时，为使 J 的一个生成族是完全割的，必须且只需它含 r 个元素。

由 Nakayama 引理，有 \(\mathrm{JM} \neq \mathrm{M}\) ，且 \(r\) 是 \(\mathbf{J}\) 的生成元的最小个数；因此推论 2 由推论 1 得出。

命题 4.--- 设 \(\rho: A \to B\) 为环同态，\(J\) 为 \(A\) 的理想，\(N\) 为 B-模。我们有等式 \(\operatorname{prof}_{\mathrm{A}}(\mathrm{J}; \mathrm{N}) = \operatorname{prof}_{\mathrm{B}}(\mathrm{JB}; \mathrm{N})\) 。

设 \(\mathbf{x} = (x_{i})_{i \in \mathbb{I}}\) 为 J 的一个生成族；族 \(\rho(\mathbf{x}) = (\rho(x_{i}))_{i \in \mathbb{I}}\) 生成 JB。按构造，复形 \(\mathbf{K}^{\bullet}(\boldsymbol{\rho}(\mathbf{x}), \mathbf{N})\) 等于 \(\mathbf{K}^{\bullet}(\mathbf{x}, \mathbf{N})\) 。因此本命题由定理 1 得出。

推论. 设 A 为局部环，a 为 A 的异于 A 的理想，M 为被 a 零化的 A-模。我们有 \(\operatorname{prof}_{\mathrm{A}}(\mathrm{M}) = \operatorname{prof}_{\mathrm{A}/\mathfrak{a}}(\mathrm{M})\) 。

设 \(\rho: A \to B\) 为环同态，\(x = (x_i)_{i \in I}\) 为 \(\Lambda\) 中元素的一个有限族，M 为 A-模。对每个整数 \(p\) ，以 \(u^p: B \otimes_\Lambda C_I^p(M) \longrightarrow C_I^p(B \otimes_\Lambda M)\) 表示把 \(b \otimes m\) 对应到交错映射 \((\alpha_1, \ldots, \alpha_p) \mapsto b \otimes m(\alpha_1, \ldots, \alpha_p)\) 的 B-线性同态。族 \((u^p)\) 定义了复形之间的一个同构

\[
u: \mathrm{B} \otimes_ {\mathrm{A}} \mathbf {K} ^ {\bullet} (\mathbf {x}, \mathrm{M}) \longrightarrow \mathbf {K} ^ {\bullet} (\mathbf {x}, \mathrm{B} \otimes_ {\mathrm{A}} \mathrm{M}).
\]

考察典范同态

\[
\gamma^ {p} (\mathrm{B}, \mathbf {K} ^ {\bullet} (\mathbf {x}, \mathrm{M})): \mathrm{B} \otimes_ {\mathrm{A}} \mathrm{H} ^ {p} (\mathbf {x}, \mathrm{M}) \longrightarrow \mathrm{H} ^ {p} (\mathrm{B} \otimes_ {\mathrm{A}} \mathbf {K} ^ {\bullet} (\mathbf {x}, \mathrm{M}))
\]

(A, X, 第 62 页)；与 \(\mathrm{H}^{p}(u)\) 复合，由它导出同态

\[
v ^ {p}: \mathrm{B} \otimes_ {\Lambda} \mathrm{H} ^ {p} (\mathbf {x}, \mathrm{M}) \longrightarrow \mathrm{H} ^ {p} (\mathbf {x}, \mathrm{B} \otimes_ {\Lambda} \mathrm{M}).
\]

引理 1.-- 若 A-模 B 平坦，则同态 \(v^p\) 对每个整数 \(p\) 都是双射。

这由 A, X, 第 66 页, 推论 2 得出。

命题 5.--- 设 A 为环，J 为 A 的有限生成理想，M 为 A-模。以 Ω 表示属于 Supp(M) 且包含 J 的 A 的极大理想所成的集。我们有

\[
\operatorname{prof} _ {A} (J; M) = \inf _ {\mathfrak {p} \in \operatorname{Spec} (A)} \operatorname{prof} _ {A _ {\mathfrak {p}}} \left(J _ {\mathfrak {p}}; M _ {\mathfrak {p}}\right) = \inf _ {m \in \Omega} \operatorname{prof} _ {\Lambda_ {m}} \left(J _ {m}; M _ {m}\right).
\]

设 \(\mathbf{x} = (x_{i})_{i \in 1}\) 为 J 的一个有限生成族。设 p 为 A 的素理想；理想 \(J_{p}\) 由像 \(x_{p}\) ——即族 \((x_{i})\) 在 \(A_{p}\) 中的像——生成。对每个 \(p \geqslant 0\) ，\(A_{p}\) -模 \(\left(\mathrm{H}^{p}(\mathbf{x},\mathrm{M})\right)_{\mathfrak{p}}\) 同构于 \(\mathrm{H}^{p}(\mathbf{x}_{\mathfrak{p}},\mathrm{M}_{\mathfrak{p}})\) （引理 1）；因此由定理 1，有 \(\operatorname{prof}_{\Lambda}(\mathrm{J};\mathrm{M}) \leqslant \inf_{\mathfrak{p} \in \operatorname{Spec}(\mathrm{A})} \operatorname{prof}_{\mathrm{A}_{\mathfrak{p}}}(\mathrm{J}_{\mathfrak{p}}; \mathrm{M}_{\mathfrak{p}}) \leqslant \inf_{\mathfrak{m} \in \Omega} \operatorname{prof}_{\mathrm{A}_{\mathfrak{m}}}(\mathrm{J}_{\mathfrak{m}}; \mathrm{M}_{\mathfrak{m}})\) 。

设 \(p\) 为一个整数，它严格小于 \(\operatorname{prof}_{\mathrm{A}_{\mathfrak{m}}}\left(\mathrm{JA}_{\mathfrak{m}};\mathrm{M}_{\mathfrak{m}}\right)\) （对每个 \(\mathfrak{m}\in\Omega\) ）。则 \(\mathrm{H}^{p}\left(\mathbf{x}_{\mathfrak{m}},\mathrm{M}_{\mathfrak{m}}\right)=0\) 对 A 的每个极大理想 \(\mathfrak{m}\) 成立：当 \(\mathfrak{m}\in\Omega\) 时这由定理 1 得出，由 \(\mathrm{M}_{\mathfrak{m}}=0\) （当 \(\mathfrak{m}\notin\operatorname{Supp}(\mathrm{M})\) 时）这一事实亦可得之，而由下述事实亦可得之：理想 \(\mathrm{JA}_{\mathfrak{m}}\) 零化 \(\mathrm{H}^{p}\left(\mathbf{x}_{\mathfrak{m}},\mathrm{M}_{\mathfrak{m}}\right)\) (A, X, 第 148 页, 推论 2)，且等于 \(\mathrm{A}_{\mathfrak{m}}\) 当 \(\mathfrak{m}\notin\mathrm{V}(\mathrm{J})\) 时。于是 \(\left(\mathrm{H}^{p}(\mathbf{x},\mathrm{M})\right)_{\mathfrak{m}}=0\) 对 A 的每个极大理想 \(\mathfrak{m}\) 成立，这蕴含 \(\mathrm{H}^{p}(\mathbf{x},\mathrm{M})=0\) (II, § 3, 第 3 节, 定理 1 的推论 2)。本命题于是由定理 1 得出。

命题 6.--- 设 A 为环，J 为 A 的有限生成理想，M 为 A-模。设 B 为环，\(\rho: A \to B\) 为使 B 成为平坦 A-模的环同态。

\begin{enumerate}
\def\labelenumi{\alph{enumi})}
\item
  有 \(a \operatorname{prof}_{\mathrm{A}}(\mathrm{J}; \mathrm{M}) \leqslant \operatorname{prof}_{\mathrm{B}}(\mathrm{JB}; \mathrm{B} \otimes_{\mathrm{A}} \mathrm{M})\) 。
\item
  再设 Supp(M) 的每个包含 J 的极大理想都属于典范映射 \(\operatorname{Spec}(B) \to \operatorname{Spec}(A)\) 的像。则 \(\operatorname{prof}_{\mathrm{A}}(\mathrm{J}; \mathrm{M}) = \operatorname{prof}_{\mathrm{B}}(\mathrm{JB}; \mathrm{B} \otimes_{\mathrm{A}} \mathrm{M})\) 。特别地，当 A-模 B 忠实平坦时即属此情形。
\end{enumerate}

断言 a) 由定理 1 与引理 1 得出。

设 \(p\) 为严格小于 \(\text{prof}_B(\text{JB};\text{B} \otimes_A \text{M})\) 的整数，并设 \(\mathfrak{m}\) 为属于 \(\text{Supp(M)} \cap \text{V(J)}\) 的 A 的极大理想。设 \(\mathbf{x}\) 为理想 J 的一个有限生成族。在 b) 的假设下，存在 B 的素理想 \(\mathfrak{n}\) 位于 \(\mathfrak{m}\) 之上，并且有一个典范同构

\[
\mathrm{B} _ {\mathfrak {n}} \otimes_ {\mathrm{A} _ {\mathfrak {m}}} \left(\mathrm{A} _ {\mathfrak {m}} \otimes_ {\Lambda} \mathrm{H} ^ {p} (\mathbf {x}, \mathrm{M})\right) \longrightarrow \mathrm{B} _ {\mathfrak {n}} \otimes_ {\mathrm{B}} \left(\mathrm{B} \otimes_ {\mathrm{A}} \mathrm{H} ^ {p} (\mathbf {x}, \mathrm{M})\right).
\]

现在 \(\mathrm{B}\otimes_{\Lambda}\mathrm{H}^{p}(\mathbf{x},\mathrm{M})\) 同构于 \(\mathrm{H}^{p}(\mathsf{p}(\mathbf{x}),\mathrm{B}\otimes_{\Lambda}\mathrm{M})\) （引理 1），故为零；此外 \(B_{n}\) 在 \(A_{m}\) 上忠实平坦 (I, § 3, 第 5 节, 命题 9 及 II, § 3, 第 4 节, 命题 14 及注)。于是有 \(A_{m}\otimes_{\Lambda}\mathrm{H}^{p}(\mathbf{x},\mathrm{M})=0\) ，从而 p \textless{} prof \(_{A_{m}}(J_{m};M_{m})\) （引理 1 与定理 1）。b) 的第一个断言于是由命题 5 得出；第二个断言由 I, § 3, 第 5 节, 命题 9 得出）。

推论.--- 设 A 为 Noether 环，J 为 A 的理想，M 为有限生成 A-模，\(\widehat{\Lambda}\) 与 \(\widehat{M}\) 为 A 与 M 关于 J-进拓扑的分离完备化。则有 \(\operatorname{prof}_{\mathbf{A}}(\mathbf{J};\mathbf{M}) = \operatorname{prof}_{\widehat{\mathbf{A}}}(\widehat{\mathbf{J}}\widehat{\mathbf{A}};\widehat{\mathbf{M}})\) 。

事实上，A-模 \(\widehat{\mathsf{A}}\) 平坦，且 \(\widehat{\mathsf{A}}\) -模 \(\widehat{\mathsf{M}}\) 同构于 \(\widehat{\mathsf{A}}\otimes_{\mathsf{A}}\mathsf{M}\) (III, § 3, 第 4 节, 定理 3)；此外，A 的每个包含 J 的极大理想都属于映射 \(\operatorname{Spec}(\widehat{\mathsf{A}})\to\operatorname{Spec}(\mathsf{A})\) 的像（同上引文, 命题 8）。

